\clearpage\pagestyle{otMain}\fontsize{11}{15}\selectfont\corrigesChapitre{1}{Nombres rationnels et fractions}
\begin{blocLivre}{cor-1-1}\begin{corrige}{1}{Une température}\demarche{Traduire chaque changement avec son signe.}
On respecte l’ordre des changements : $-6+11=5$, puis $5-4=\res{1\,^\circ\mathrm C}$. Une augmentation se traduit par une addition, même si la température initiale est négative.
\end{corrige}\end{blocLivre}
\begin{blocLivre}{cor-1-2}\begin{corrige}{2}{Respecter les priorités}\demarche{Calculer les parenthèses avant les produits.}
Dans $A$, la parenthèse vaut 7, puis $3\times7=21$ : $\res{A=-3}$. Dans $B$, les deux parenthèses donnent 15 et 7 : $\res{B=105}$. Les parenthèses ont changé les facteurs du produit.
\end{corrige}\end{blocLivre}
\begin{blocLivre}{cor-1-3}\begin{corrige}{3}{Deux écritures, un nombre}\demarche{Appliquer le même facteur aux deux termes.}
On multiplie 8 par 5, donc aussi 5 par 5 : $\res{25/40}$. On divise 42 par 21, donc aussi 63 par 21 : $\res{2/3}$. La même opération sur les deux termes conserve le quotient.
\end{corrige}\end{blocLivre}
\begin{blocLivre}{cor-1-4}\begin{corrige}{4}{Une somme exacte}\demarche{Choisir un dénominateur commun adapté.}
Un dénominateur commun est 36 : $7/12=21/36$ et $5/18=10/36$. La somme est $\res{31/36}$. Le numérateur 31 n’a aucun diviseur commun avec 36 autre que 1.
\end{corrige}\end{blocLivre}
\begin{blocLivre}{cor-1-5}\begin{corrige}{5}{Produit et quotient}\demarche{Simplifier les facteurs et inverser le diviseur.}
Dans le produit, on simplifie 7 avec 14 et 9 avec 15 : $\res{-3/10}$. Dans le quotient, on multiplie par l’inverse : $(3/5)\times(-10/9)=\res{-2/3}$. Les deux résultats sont négatifs.
\end{corrige}\end{blocLivre}
\begin{blocLivre}{cor-1-6}\begin{corrige}{6}{Comparer sans arrondir}\demarche{Ramener les fractions à un dénominateur positif commun.}
Avec le dénominateur 12, les numérateurs sont $-9,-10,-7$. Ainsi $\res{-5/6<-3/4<-7/12}$. Parmi des négatifs, le nombre le plus éloigné de zéro est le plus petit.
\end{corrige}\end{blocLivre}
\begin{blocLivre}{cor-1-7}\begin{corrige}{7}{Une expression à étages}\demarche{Séparer le numérateur et le dénominateur avant de diviser.}
Le numérateur vaut $4/6+1/6=5/6$ ; le dénominateur vaut $5/4-2/4=3/4$, non nul. Donc $C=(5/6)\times(4/3)=\res{10/9}$.
\end{corrige}\end{blocLivre}
\begin{blocLivre}{cor-1-8}\begin{corrige}{8}{Une réduction trompeuse}\demarche{Revenir au sens de la barre de fraction.}
La division porte sur toute la somme : $(8+4)/4=12/4=\res3$. On peut aussi distribuer la division : $8/4+4/4=2+1$. Simplifier seulement le terme 4 a laissé le 8 sans division.
\end{corrige}\end{blocLivre}
\begin{blocLivre}{cor-1-9}\begin{corrige}{9}{Une recette}\demarche{Passer par une personne ou par le coefficient de proportionnalité.}
Pour une personne, il faut $(3/4)/6=1/8$ L. Pour dix : $10/8=\res{5/4\text{ L}}$, soit $\res{1{,}25\text{ L}}$. Le facteur global $10/6$ donne le même résultat.
\end{corrige}\end{blocLivre}
\begin{blocLivre}{cor-1-10}\begin{corrige}{10}{Exact ou approché ?}\demarche{Conserver la fraction exacte avant d’arrondir.}
On écrit $2=14/7$ : le résultat exact est $\res{9/7}$. Sa valeur décimale commence par $1{,}2857\ldots$ ; le troisième chiffre après la virgule impose l’arrondi $\res{1{,}29}$. On utilise le signe $\approx$ pour cet arrondi.
\end{corrige}\end{blocLivre}
\begin{blocLivre}{cor-1-11}\begin{corrige}{11}{Deux étapes de lecture}\demarche{Identifier le tout auquel se rapporte chaque fraction.}
Après lundi, il reste $3/5$. Mardi, il lit $(1/3)\times(3/5)=1/5$ du livre entier. Il reste donc $3/5-1/5=\res{2/5}$. Cela représente $270\times2/5=\res{108\text{ pages}}$. Additionner $2/5$ et $1/3$ aurait confondu les références.
\end{corrige}\end{blocLivre}
\begin{blocLivre}{cor-1-12}\begin{corrige}{12}{Une réserve d’eau}\demarche{Exprimer les volumes comme fractions d’une même capacité.}
Après consommation, la part remplie est $(7/8)(1-2/7)=5/8$. Les 30 L ajoutés représentent $3/4-5/8=1/8$ de la capacité. La capacité vaut donc $30\times8=\res{240\text{ L}}$. Contrôle : 210 L au départ, puis 150 L, puis 180 L, soit $3/4$ de 240 L.
\end{corrige}\end{blocLivre}
\begin{blocLivre}{cor-1-13}\begin{corrige}{13}{Partager un trajet}\demarche{Calculer la part restante avant de retrouver le tout.}
La part déjà parcourue vaut $9/24+10/24=19/24$. Il reste $5/24$, qui correspond à 15 km. Une part de $1/24$ vaut 3 km, donc le trajet mesure $\res{72\text{ km}}$. On vérifie $27+30+15=72$.
\end{corrige}\end{blocLivre}
\begin{blocLivre}{cor-1-14}\begin{corrige}{14}{Des votes et des absents}\demarche{Multiplier les proportions lorsque les ensembles sont emboîtés.}
La part de toute la classe qui vote A est $(5/6)(3/5)=1/2$. Les 15 voix représentent donc la moitié de la classe : $\res{30\text{ élèves}}$. Il y a $30/6=\res{5\text{ absents}}$ et 25 présents ; $3/5$ de 25 vaut bien 15.
\end{corrige}\end{blocLivre}
\begin{blocLivre}{cor-1-15}\begin{corrige}{15}{Trouver toutes les possibilités}\demarche{Transformer les comparaisons en un encadrement entier.}
Les trois fractions se comparent avec le dénominateur 12 : $2/3=8/12$ et $5/4=15/12$. Il faut donc $8<n<15$. Les solutions sont $\res{9,10,11,12,13,14}$. Les deux bornes sont exclues à cause du mot « strictement ».
\end{corrige}\end{blocLivre}
